\documentclass[10.5pt,a4paper]{article}
\usepackage[margin=20mm]{geometry}
\usepackage{amsmath,amssymb}
\usepackage{booktabs}
\usepackage{graphicx}
\usepackage{microtype}
\usepackage[hidelinks]{hyperref}
\setlength{\parindent}{0pt}
\setlength{\parskip}{0.5em}
\graphicspath{{out/experiments/}}
\newcommand{\Prob}{\mathbb{P}}

\title{Experiments: Robustness and Value of Two and Three Free Intervals}
\author{}
\date{}

\begin{document}
\maketitle

\section*{Setup}
Model and notation as in \texttt{allocation\_math.pdf}. Defaults: \(N=40\), \(T=1000\) A,
\((I_S,I_M,I_L)=(20,40,100)\) A, \(P_j=VI_j\) with \(V=240\) V, mix \((p_S,p_M,p_L)=(0.69,0.30,0.01)\),
uptake \(u=10\%\), All-On, 2025 VIC1 data in market time (UTC+10). Violation uses each circuit's
lowest-peak allocation; soak uses each circuit's best safe allocation (safest end of the Pareto
front). Code: \texttt{play/experiments/}, run with \texttt{python3 -m play.experiments}.

\begin{center}\small
\begin{tabular}{lll}
\toprule
Design & Intervals & Non-overlapping intervals\\
\midrule
1\(\times\)3h & 11--14 & 1\\
2\(\times\)3h & 09--12, 12--15 & 2\\
3\(\times\)3h overlap & 08--11, 09--12, 12--15 & 2\\
3\(\times\)2h & 09--11, 11--13, 13--15 & 3\\
3\(\times\)3h no overlap & 08--11, 11--14, 14--17 & 3\\
\bottomrule
\end{tabular}
\end{center}

\section{E1 On-Off usage}
Each medium/large customer's extra load is on for its whole interval with probability \(q\)
(base load always on); the operator allocates with the All-On rule. For interval \(k\),
\(Z_k=(I_M-I_S)Y^M_k+(I_L-I_S)Y^L_k\) with \(Y^M_k\sim\mathrm{Bin}(m_k,q)\),
\(Y^L_k\sim\mathrm{Bin}(l_k,q)\), independent, and
\begin{equation}
  \Prob(\text{no overload}\mid\boldsymbol n)=\Prob\Bigl(NI_S+\sum_{k\in C(g)}Z_k\le T\ \ \forall g\Bigr),
\end{equation}
computed exactly over the joint distribution of \((Z_1,\dots,Z_K)\); \(q=1\) is All-On.

\begin{center}\small
\begin{tabular}{lrrrrr}
\toprule
\(\Prob(\text{violation})\) & \(q=0.1\) & 0.25 & 0.5 & 0.75 & 1\\
\midrule
1\(\times\)3h & 0.009\% & 0.495\% & 10.28\% & 42.98\% & 79.59\%\\
2\(\times\)3h, 3\(\times\)3h overlap & 0.001\% & 0.011\% & 0.152\% & 1.065\% & 4.195\%\\
3\(\times\)2h, 3\(\times\)3h no overlap & 0.00001\% & 0.00004\% & 0.0006\% & 0.004\% & 0.014\%\\
\bottomrule
\end{tabular}
\end{center}
The ranking holds for every \(q\), and the gain grows as usage falls: 1 vs 2 intervals is
19\(\times\) at \(q=1\) and 68\(\times\) at \(q=0.5\).

\section{E2 Customer mix}
All-On violation over \(p_M\in[0,0.5]\), \(p_L\in[0,0.1]\) (\(p_S=1-p_M-p_L\)), Fig.~\ref{fig:mix}.
Two intervals keep violation below 5\% while \(p_M\le0.30\) and \(p_L\le0.01\), or \(p_M\le0.20\) and
\(p_L\le0.02\); three non-overlapping 2 h intervals stay below 1\% up to \(p_L=0.03,\ p_M\le0.30\).
\begin{figure}[h]\centering\includegraphics[width=\textwidth]{e02_mix.png}
\caption{Violation (\%) by customer mix.}\label{fig:mix}\end{figure}

\section{E3 Circuit size}
\(N\in\{20,40,60,80\}\) with \(T=25N\) A. With average All-On current per customer
\(\mu=\sum_jp_jI_j=26.8\) A and \(K\) balanced non-overlapping intervals, the busiest interval
carries about \(N\bigl(I_S+(\mu-I_S)/K\bigr)\): \(26.8N\) for \(K=1\), \(23.4N\) for \(K=2\),
\(22.3N\) for \(K=3\). Against \(25N\), the law of large numbers drives violation to 1 for one interval and to 0
for two or more as \(N\) grows.
\begin{center}\small
\begin{tabular}{lrrrr}
\toprule
\(\Prob(\text{violation})\) & \(N=20\) & 40 & 60 & 80\\
\midrule
1 interval & 66.7\% & 79.6\% & 86.5\% & 90.7\%\\
2 non-overlapping & 8.8\% & 4.2\% & 2.1\% & 1.0\%\\
3 non-overlapping & 0.37\% & 0.014\% & 0.0005\% & 0.00002\%\\
\bottomrule
\end{tabular}
\end{center}

\section{E4/E5 Same daily energy}
The window 09--15 is split into \(K\) equal slots of length \(L=6/K\) h; every medium/large
customer needs \(H=2\) h of full-power use. \textbf{E4 inflexible}: full power at an uncontrolled
time, so overload uses the full current over the slot (worst case). \textbf{E5 flexible} (smart
charging): the energy is spread over the slot, so the extra current is \((I_j-I_S)H/L\). In both,
soak uses the expected load \((P_j-P_S)H/L\) over the slot.
\begin{center}\small
\begin{tabular}{lrrrrr}
\toprule
Slots & E4 violation & E5 violation & soak \(u=1\%\) & 10\% & 50\%\\
\midrule
1\(\times\)6h & 79.59\% & 0.014\% & 92.5 & 62.8 & 24.2\\
2\(\times\)3h & 4.195\% & 0.014\% & 92.6 & 62.8 & 24.2\\
3\(\times\)2h & 0.014\% & 0.014\% & 92.8 & 62.8 & 24.2\\
\bottomrule
\end{tabular}
\end{center}
With equal energy, soak is the same for every \(K\). Splitting only matters for inflexible
loads: multiple intervals do the same job as spreading the load, so they substitute for load control.

\section{E6 Daylight saving}
Interval times as local clock time: on daylight-saving days local \(t\) = market \(t+1\) h, so
\(x^{\text{local}}_d(t)=x_d(t-1\,\text{h})\). Best 3 h designs inside 08--15 are unchanged
(2: 09--12, 12--15; 3: 08--11, 09--12, 12--15). Soak falls by 2--8\% (e.g.\ 80.5 \(\to\) 79.1 kWh/day for
2 intervals at 10\%), because in summer these clock times sit an hour earlier than the
curtailment peak. The gain of 3 over 2 drops to 0.4--5\%.

\section{E7 Rooftop-solar timing}
Soak against the VIC1 rooftop-solar shape, scaled to the same annual energy as the curtailment so only timing
changes. The ranking is unchanged and the best designs barely move. Rooftop output is
concentrated at midday, so at 10\% uptake soak rises (2 intervals: 80.5 \(\to\) 120.3 kWh/day).
The gain of 3 over 2 is 0.0\%, 2.6\% and 6.4\% at 1\%, 10\% and 50\% uptake (curtailment: 0.4\%, 1.5\%, 7.9\%).

\section{E8 Evening peak}
Third objective: added load in each day's top 6 demand hours,
\(\sum_t \mathbb E[D(t)]\,f(t)\,\Delta t\), where \(f(t)\) is the share of days in which slot \(t\) is in that day's top
6 demand hours (morning 06--09 on winter days, 16--22 most days). 3 h designs inside 08--17, \(u=10\%\).
Selected Pareto-optimal designs:
\begin{center}\small
\begin{tabular}{lrrr}
\toprule
Design & violation & soak & peak load (kWh/day)\\
\midrule
11--14 & 79.6\% & 58.0 & 2.5\\
09--12, 12--15 & 4.2\% & 80.5 & 4.9\\
09--12, 12--15, 13--16 & 4.2\% & 82.7 & 5.7\\
08--11, 10--13, 13--16 & 4.2\% & 84.0 & 9.3\\
08--11, 11--14, 14--17 & 0.014\% & 84.8 & 15.0\\
\bottomrule
\end{tabular}
\end{center}
The safest design, three non-overlapping intervals to 17:00, puts 3\(\times\) more load into peak
hours than the best 2-interval design. Ending by 15:00--16:00 keeps peak load low at 4.2\% violation.

\section{E9 Grid ramp}
State added load \(D_{\text{state}}(t)=C\,\mathbb E[D(t)]\) (MW), \(C=uN_{\mathrm{VIC}}/N=7{,}250\);
the largest 5-minute change is the amount switching at an interval edge.
\begin{center}\small
\begin{tabular}{lrr}
\toprule
Design & largest step (MW) & peak added load (MW)\\
\midrule
08--11, 11--14 & 275 & 275\\
09--12, 12--15 & 248 & 248\\
08--11, 09--12, 12--15 & 235 & 239\\
08--11, 10--13, 12--15 & 198 & 317\\
09--11, 11--13, 13--15 & 158 & 175\\
08--11, 11--14, 14--17 & 141 & 212\\
\bottomrule
\end{tabular}
\end{center}
More start times split the switching into smaller steps, which gives smaller ramps. Overlap can
raise the peak added load, because the allocation here maximises soak, not smoothness.

\section{E10 Value}
Upgrades avoided against one interval: \((P_1-P_K)\,C\). At \(u=10\%\): 5{,}466 circuits (2 intervals)
and 5{,}769 (3 non-overlapping); at \(u=100\%\): 54{,}662 and 57{,}693. At an \emph{illustrative}
AUD 100{,}000 per upgrade (placeholder; replace with DNSP cost data) that is AUD 0.55 bn at 10\% uptake.

\section*{Conclusions}
\begin{itemize}\itemsep0pt
\item Two non-overlapping intervals cut violation 19--68\(\times\) against one interval. This holds for usage
  probability (E1), most realistic mixes (E2), every circuit size (E3, growing with \(N\)), daylight
  saving (E6) and solar timing (E7).
\item Three intervals improve safety only if they do not all overlap. Three non-overlapping intervals give 0.014\%,
  but either need a window wider than two can cover (E8: more evening-peak load) or shorter intervals.
\item With equal daily energy (E4/E5), soak does not depend on \(K\). The value of extra intervals is
  safety for inflexible loads, plus smaller grid ramps (E9).
\item The case rests on mix and limit data (E2, E3) and upgrade cost (E10), which need real DNSP inputs.
\end{itemize}

\end{document}
